% ==============================================================================
% 08_ui_ux.tex — Phần 8. Thiết kế giao diện và trải nghiệm (UI/UX)
% Phụ trách: CHƯA GIAO (Tạm thời không giao cho Thành viên 3)
% Ngân sách trang: 4 trang
% ==============================================================================

\pttkPart{Thiết kế giao diện và trải nghiệm (UI/UX)}

\todosection{CHƯA GIAO}{Phần 8 tạm thời chưa phân công (không bàn giao cho Thành viên 3). Tiến và Xuân sẽ thống nhất phân công sau. Ngân sách mục tiêu: 4 trang.}

\subsection{Chân dung người dùng (Personas)}
% Bảng 3 Personas: Solo Dev, Tech Lead nhóm nhỏ, Lập trình viên Frontend/Backend.

\subsection{Hành trình người dùng (User Journey)}
% Bảng hành trình nhận Task -> Code & Push Commit -> CI Fail được AI tóm tắt -> Merge PR tự động Done.

\subsection{Danh sách màn hình}
% Bảng SC-01..SC-10 ánh xạ theo glossary\_lock.md và các trang trong frontend/src/pages/.

\subsection{Wireframe màn hình Kanban và AI ChatPanel (SC-07, SC-09)}
\todofig{Hình 9. Wireframe / Giao diện màn hình Kanban kèm Git/CI Badges và AI ChatPanel — diagrams/out/hinh\_09\_wireframe\_kanban.png}

\subsection{Hệ thống thiết kế (Design System)}
% Bảng màu TailwindCSS v4, Typography, hệ lưới, component tái sử dụng (Task Card, Git/CI Badge, Action Card, ProblemDetail Toast).

\subsection{Khả năng truy cập (WCAG 2.1 AA)}
% Độ tương phản >= 4.5:1, không chỉ dùng màu sắc, thao tác bàn phím/menu thay thế kéo thả Kanban.
